\documentclass[aspectratio=169]{beamer}

\usetheme{Madrid}
\usepackage[spanish]{babel}
\usepackage[utf8]{inputenc}
\usepackage{textcomp}
\usepackage[T1]{fontenc}
\usepackage{lmodern}
\usepackage{graphicx}
\usepackage{booktabs}
\usepackage{verbatim}
\usepackage{amssymb}
\usepackage{pifont}
\usepackage{moreverb}
\usepackage[most]{tcolorbox}
\usepackage{tikz}
\usetikzlibrary{shapes, arrows, positioning}

% Configuración de TikZ
\tikzstyle{block} = [
  rectangle, draw, rounded corners,
  text centered, minimum height=1cm,
  minimum width=2.5cm, font=\small, fill=gray!5
]
\tikzstyle{arrow} = [->, thick]

% Configuración de tcolorbox
\tcbset{
  colback=gray!5, colframe=gray!50,
  boxrule=0.4pt, sharp corners,
}

\title[BookMate]{BookMate — Sistema de Recomendación de Libros Académicos}
\subtitle{Charter del Proyecto y Modelo de Negocio}
\author[GRUPO 6]{Delgado R., G. \and Osorio M., A. \and Rojas A., J. \and Torres R., J. \and Valverde G., Y. \and Villanueva A., F.}
\institute[FC-UNI]{UNIVERSIDAD NACIONAL DE INGENIERÍA}
\titlegraphic{\includegraphics[height=1.4cm]{example-image}}
% NOTA: Reemplazar 'example-image' con 'Imagenes/logo_uni.png' cuando esté disponible

\date{Octubre 2025}

\begin{document}

% Portada
\begin{frame}
\titlepage
\end{frame}

% Agenda
\begin{frame}
\frametitle{Agenda}
\tableofcontents
\end{frame}

\section{Introducción al Proyecto}

\begin{frame}
\frametitle{¿Qué es BookMate?}

\begin{columns}
\begin{column}{0.5\textwidth}
\textbf{Descripción:}
\begin{itemize}
\item Sistema de gestión de libros académicos
\item Recomendaciones inteligentes con IA
\item Análisis semántico de contenido
\item Interfaz web moderna
\end{itemize}

\vspace{0.5cm}

\textbf{Objetivo Principal:}
\begin{tcolorbox}[colback=blue!5, colframe=blue!50]
Facilitar el descubrimiento de libros académicos mediante recomendaciones basadas en similitud semántica.
\end{tcolorbox}
\end{column}

\begin{column}{0.5\textwidth}
\begin{center}
\includegraphics[width=\textwidth]{example-image-a}
% NOTA: Reemplazar con captura de pantalla del sistema
% o diagrama conceptual de BookMate
\end{center}
\end{column}
\end{columns}

\end{frame}

\begin{frame}
\frametitle{Problema a Resolver}

\textbf{Desafíos actuales en búsqueda de libros académicos:}

\begin{enumerate}
\item \textbf{Sobrecarga de información}
  \begin{itemize}
  \item Miles de libros disponibles
  \item Difícil identificar los más relevantes
  \end{itemize}

\item \textbf{Búsquedas limitadas}
  \begin{itemize}
  \item Solo coincidencias exactas de texto
  \item No comprenden el contenido real
  \end{itemize}

\item \textbf{Recomendaciones superficiales}
  \begin{itemize}
  \item Basadas únicamente en género o autor
  \item No analizan similitud de contenido
  \end{itemize}
\end{enumerate}

\vspace{0.3cm}

\begin{tcolorbox}[colback=red!5, colframe=red!50]
\textbf{Resultado:} Tiempo desperdiciado y libros relevantes no descubiertos
\end{tcolorbox}

\end{frame}

\begin{frame}
\frametitle{Solución Propuesta}

\begin{columns}
\begin{column}{0.4\textwidth}
\textbf{Recomendaciones con IA:}
\begin{itemize}
\item Análisis de sinopsis con NLP
\item Embeddings semánticos
\item Similitud de coseno
\item Recomendaciones relevantes
\end{itemize}

\vspace{0.3cm}

\textbf{Sistema Híbrido:}
\begin{itemize}
\item Recomendaciones heurísticas
\item Fallback automático
\item Alta disponibilidad
\end{itemize}
\end{column}

\begin{column}{0.6\textwidth}
\begin{center}
\begin{tikzpicture}[node distance=2cm]
\node[block] (usuario) {Usuario};
\node[block, right=of usuario] (sistema) {Sistema BookMate};
\node[block, above right=1cm and 0.5cm of sistema] (ia) {Servicio IA};
\node[block, below right=1cm and 0.5cm of sistema] (heur) {Heurísticas};
\draw[arrow] (usuario) -- (sistema);
\draw[arrow] (sistema) -- node[above, sloped] {primario} (ia);
\draw[arrow] (sistema) -- node[below, sloped] {fallback} (heur);
\end{tikzpicture}
% NOTA: Este diagrama representa la arquitectura de recomendaciones
\end{center}
\end{column}
\end{columns}

\end{frame}

\section{Charter del Proyecto}

\begin{frame}
\frametitle{Propósito y Visión}

\textbf{Propósito del Proyecto:}
\begin{itemize}
\item Desarrollar un sistema de recomendación de libros académicos
\item Utilizar inteligencia artificial para análisis semántico
\item Proporcionar recomendaciones personalizadas
\item Facilitar la gestión eficiente del catálogo
\end{itemize}

\vspace{0.5cm}

\textbf{Visión:}
\begin{tcolorbox}[colback=green!5, colframe=green!50]
BookMate será la plataforma de referencia para el descubrimiento de libros académicos, distinguiéndose por su capacidad de comprender el contenido semántico y proporcionar recomendaciones altamente relevantes.
\end{tcolorbox}

\end{frame}

\begin{frame}
\frametitle{Alcance del Proyecto}

\begin{columns}
\begin{column}{0.5\textwidth}
\textbf{Funcionalidades Incluidas:}
\begin{itemize}
\item[\checkmark] Gestión de catálogo (CRUD)
\item[\checkmark] Búsqueda avanzada
\item[\checkmark] Recomendaciones con IA
\item[\checkmark] Recomendaciones heurísticas
\item[\checkmark] Interfaz web responsive
\item[\checkmark] Panel de administración
\end{itemize}
\end{column}

\begin{column}{0.5\textwidth}
\textbf{Fuera de Alcance:}
\begin{itemize}
\item[\ding{55}] Autenticación completa
\item[\ding{55}] Sistema de compras
\item[\ding{55}] APIs externas de libros
\item[\ding{55}] Notificaciones por email
\item[\ding{55}] Aplicación móvil nativa
\item[\ding{55}] Sistema de reseñas
\end{itemize}
\end{column}
\end{columns}

\end{frame}

\begin{frame}
\frametitle{Stakeholders}

\begin{table}
\centering
\small
\begin{tabular}{p{3cm}p{3cm}p{5cm}}
\toprule
\textbf{Stakeholder} & \textbf{Rol} & \textbf{Expectativas} \\
\midrule
Equipo de Desarrollo & Desarrolladores & Producto funcional y bien documentado \\
Profesor del Curso & Cliente/Evaluador & Cumplimiento de objetivos académicos \\
Usuarios Finales & Usuarios del sistema & Sistema intuitivo y recomendaciones relevantes \\
Administradores & Gestores & Herramientas eficientes para gestión \\
\bottomrule
\end{tabular}
\end{table}

\end{frame}

\begin{frame}
\frametitle{Requisitos de Alto Nivel}

\textbf{Requisitos Funcionales Principales:}

\begin{table}
\centering
\small
\begin{tabular}{p{1.5cm}p{7cm}p{2cm}}
\toprule
\textbf{ID} & \textbf{Requisito} & \textbf{Prioridad} \\
\midrule
RF-01 & Gestionar libros (CRUD completo) & Alta \\
RF-02 & Gestionar autores (CRUD completo) & Alta \\
RF-03 & Búsqueda por título, autor, género & Alta \\
RF-04 & Generar recomendaciones personalizadas & \textbf{Muy Alta} \\
RF-05 & Mostrar detalles completos de libros & Media \\
RF-06 & Filtrar libros por precio & Media \\
\bottomrule
\end{tabular}
\end{table}

\vspace{0.3cm}

\begin{tcolorbox}[colback=yellow!5, colframe=yellow!50]
\textbf{RF-04} es el requisito más crítico del proyecto
\end{tcolorbox}

\end{frame}

\begin{frame}
\frametitle{Requisitos No Funcionales}

\begin{table}
\centering
\small
\begin{tabular}{p{4cm}p{4.5cm}p{2cm}}
\toprule
\textbf{Categoría} & \textbf{Requisito} & \textbf{Objetivo} \\
\midrule
Rendimiento & Búsquedas & <1s (p95) \\
Rendimiento & Recomendaciones IA & <3s (p95) \\
Usabilidad & Interfaz intuitiva & <5 clicks \\
Usabilidad & Responsive design & Todos los dispositivos \\
Fiabilidad & Disponibilidad & 95\% \\
Seguridad & HTTPS & Producción \\
Mantenibilidad & Cobertura tests & >80\% \\
Escalabilidad & Usuarios concurrentes & 100 \\
\bottomrule
\end{tabular}
\end{table}

\end{frame}

\begin{frame}
\frametitle{Restricciones y Supuestos}

\begin{columns}
\begin{column}{0.5\textwidth}
\textbf{Restricciones:}
\begin{itemize}
\item \textbf{Tiempo:} 10 semanas (S7-S16)
\item \textbf{Equipo:} 6 integrantes
\item \textbf{Recursos:} Entorno local
\item \textbf{Fases:}
  \begin{itemize}
  \item Análisis: S7-S11
  \item Diseño: S13-S16
  \end{itemize}
\end{itemize}
\end{column}

\begin{column}{0.5\textwidth}
\textbf{Supuestos:}
\begin{itemize}
\item Disponibilidad del equipo: 8-10h/semana
\item Acceso a herramientas de desarrollo
\item Uso de modelo preentrenado (no entrenamiento)
\item Datos ingresados manualmente
\item Prioridad: prototipo funcional
\end{itemize}
\end{column}
\end{columns}

\end{frame}

\begin{frame}
\frametitle{Riesgos Principales}

\begin{table}
\centering
\tiny
\begin{tabular}{p{3.5cm}p{1.5cm}p{1.5cm}p{4.5cm}}
\toprule
\textbf{Riesgo} & \textbf{Prob.} & \textbf{Impacto} & \textbf{Mitigación} \\
\midrule
Complejidad servicio IA & Media & Alto & Fallback a heurísticas; modelo preentrenado \\
Integración microservicios & Media & Medio & Documentar APIs; Docker Compose \\
Tiempo insuficiente & Alta & Alto & Priorizar features core \\
Problemas de rendimiento IA & Media & Medio & Caché de embeddings \\
Disponibilidad del equipo & Media & Medio & Reuniones semanales \\
\bottomrule
\end{tabular}
\end{table}

\end{frame}

\section{Visión del Producto}

\begin{frame}
\frametitle{Declaración de Visión}

\begin{tcolorbox}[colback=blue!5, colframe=blue!50, title=Visión de BookMate]
\large
BookMate será la plataforma líder en gestión y descubrimiento inteligente de libros académicos, transformando la manera en que estudiantes y profesionales encuentran contenido relevante mediante el poder de la inteligencia artificial y el análisis semántico.
\end{tcolorbox}

\vspace{0.5cm}

\textbf{Valor Diferencial:}
\begin{itemize}
\item Comprensión del contenido real, no solo metadatos
\item Descubrimiento serendípico de libros relacionados
\item Sistema confiable con fallback automático
\item Enfoque específico en literatura académica
\end{itemize}

\end{frame}

\begin{frame}
\frametitle{Usuarios Objetivo}

\begin{columns}
\begin{column}{0.33\textwidth}
\textbf{Estudiantes}
\begin{itemize}
\item Buscan bibliografía
\item Para tesis/proyectos
\item Aprendizaje autodirigido
\end{itemize}

\vspace{0.3cm}

\textbf{Beneficio:}
Descubrir libros relevantes que no conocían
\end{column}

\begin{column}{0.33\textwidth}
\textbf{Profesores}
\begin{itemize}
\item Bibliografía para cursos
\item Material de investigación
\item Literatura relacionada
\end{itemize}

\vspace{0.3cm}

\textbf{Beneficio:}
Identificar rápidamente material complementario
\end{column}

\begin{column}{0.33\textwidth}
\textbf{Bibliotecarios}
\begin{itemize}
\item Gestionar catálogos
\item Ayudar a usuarios
\item Mantener actualizado
\end{itemize}

\vspace{0.3cm}

\textbf{Beneficio:}
Herramientas eficientes de administración
\end{column}
\end{columns}

\end{frame}

\begin{frame}
\frametitle{Escenario de Uso: Estudiante}

\textbf{Contexto:} María investiga sobre "arquitectura de microservicios" para su tesis

\vspace{0.3cm}

\begin{enumerate}
\item Busca un libro conocido sobre microservicios
\item Visualiza detalles del libro
\item Solicita recomendaciones con IA
\item BookMate analiza la sinopsis con NLP
\item Calcula similitud semántica
\item Sugiere 6 libros relacionados:
  \begin{itemize}
  \item Patrones de diseño distribuido
  \item Docker y contenedores
  \item Event-driven architecture
  \item Domain-driven design
  \end{itemize}
\item María descubre 4 libros relevantes que no conocía
\end{enumerate}

\vspace{0.3cm}

\begin{tcolorbox}[colback=green!5, colframe=green!50]
\textbf{Resultado:} 70\% de reducción en tiempo de búsqueda de bibliografía
\end{tcolorbox}

\end{frame}

\section{Modelo de Negocio (BPMN)}

\begin{frame}
\frametitle{Proceso de Recomendación - Vista General}

\begin{center}
\includegraphics[width=0.8\textwidth]{example-image-b}
% NOTA: Reemplazar con el diagrama BPMN generado desde bpmn.puml
% Mostrar el proceso simplificado de recomendación:
% Usuario → Sistema → IA/Heurísticas → Resultados
\end{center}

\textbf{Flujo principal:}
\begin{enumerate}
\item Usuario explora catálogo y selecciona libro
\item Sistema solicita recomendaciones al servicio IA
\item Servicio IA genera embeddings y calcula similitud
\item Sistema retorna top 6 libros más similares
\item Usuario visualiza recomendaciones
\end{enumerate}

\end{frame}

\begin{frame}
\frametitle{Proceso de Recomendación con IA}

\begin{center}
\includegraphics[width=0.85\textwidth]{example-image-c}
% NOTA: Reemplazar con el diagrama BPMN de interacción con IA
% Mostrar el flujo detallado:
% 1. Sistema recibe solicitud
% 2. Consulta datos del libro
% 3. Envía a servicio IA
% 4. IA genera embeddings
% 5. Calcula similitud de coseno
% 6. Retorna recomendaciones
% 7. Sistema consulta detalles
% 8. Retorna al usuario
\end{center}

\textbf{Fallback:} Si IA no disponible → Recomendaciones heurísticas

\end{frame}

\begin{frame}
\frametitle{Proceso de Administración del Catálogo}

\begin{center}
\includegraphics[width=0.7\textwidth]{example-image}
% NOTA: Reemplazar con el diagrama BPMN de administración
% Mostrar las operaciones CRUD:
% - Agregar libro
% - Editar libro
% - Eliminar libro
% - Listar libros
\end{center}

\textbf{Operaciones disponibles:}
\begin{itemize}
\item Agregar libro (con validación y generación de embedding)
\item Editar libro (regenera embedding si sinopsis cambió)
\item Eliminar libro (elimina también el embedding)
\item Listar todos los libros del catálogo
\end{itemize}

\end{frame}

\begin{frame}
\frametitle{Proceso de Búsqueda}

\textbf{Flujo de búsqueda:}

\begin{enumerate}
\item Usuario ingresa término de búsqueda
\item Sistema busca coincidencias en:
  \begin{itemize}
  \item Título
  \item Autor
  \item Género
  \end{itemize}
\item Aplica filtros opcionales:
  \begin{itemize}
  \item Precio (min/max)
  \item Género
  \item Año de publicación
  \item Rating
  \end{itemize}
\item Ordena por relevancia
\item Retorna resultados
\end{enumerate}

\vspace{0.3cm}

\begin{tcolorbox}[colback=blue!5, colframe=blue!50]
\textbf{Objetivo:} Búsquedas en <1 segundo (p95)
\end{tcolorbox}

\end{frame}

\section{Hitos del Proyecto}

\begin{frame}
\frametitle{Cronograma de Entregables}

\begin{table}
\centering
\small
\begin{tabular}{p{1.5cm}p{3.5cm}p{6cm}}
\toprule
\textbf{Semana} & \textbf{Hito} & \textbf{Entregable Principal} \\
\midrule
\textbf{7} & Inicio del proyecto & Charter, Visión, Glosario, BPMN \\
\textbf{8} & Análisis de requisitos & SRS, Casos de Uso, Prototipo UC \\
\textbf{11} & Análisis conceptual & Modelo de Análisis, C4 C1-C2, Robustez, CRC \\
\textbf{13} & Diseño arquitectónico & Arquitectura, C4 C3-C4, Plan de desarrollo \\
\textbf{15} & Diseño detallado & Clases, Secuencia, Despliegue \\
\textbf{16} & Entrega final & Demo, Planes de pruebas/despliegue/configuración \\
\bottomrule
\end{tabular}
\end{table}

\vspace{0.3cm}

\textbf{Nota:} Semanas 7-11 = Análisis (sin tecnologías) | Semanas 13-16 = Diseño (con implementación)

\end{frame}

\begin{frame}
\frametitle{Criterios de Éxito}

El proyecto será exitoso si cumple:

\begin{enumerate}
\item \checkmark \textbf{Funcionalidad core:} CRUD + Recomendaciones IA funcionando
\item \checkmark \textbf{RNF cumplidos:} Búsquedas <1s, Recomendaciones <3s
\item \checkmark \textbf{Documentación completa:} Todos los entregables presentados
\item \checkmark \textbf{Calidad de código:} Cobertura >80\%, código limpio
\item \checkmark \textbf{Demo funcional:} Demostración exitosa en semana 16
\end{enumerate}

\vspace{0.5cm}

\begin{tcolorbox}[colback=green!5, colframe=green!50, title=Meta Principal]
Sistema de recomendaciones con >70\% de precisión y respuesta en <3 segundos
\end{tcolorbox}

\end{frame}

\section{Conclusiones}

\begin{frame}
\frametitle{Resumen}

\textbf{BookMate en síntesis:}

\begin{itemize}
\item Sistema de gestión y recomendación de libros académicos
\item Utiliza IA y NLP para análisis semántico
\item Sistema híbrido (IA + heurísticas) para confiabilidad
\item Interfaz web responsive e intuitiva
\item Proyecto académico de 10 semanas (S7-S16)
\end{itemize}

\vspace{0.5cm}

\textbf{Valor principal:}
\begin{tcolorbox}[colback=blue!5, colframe=blue!50]
Transformar la búsqueda pasiva de libros en un descubrimiento activo e inteligente mediante IA
\end{tcolorbox}

\end{frame}

\begin{frame}
\frametitle{Próximos Pasos}

\textbf{Semana 8 - Análisis de Requisitos:}

\begin{itemize}
\item Elaborar Software Requirements Specification (SRS) completo
\item Definir todos los casos de uso (UC-01 a UC-15)
\item Crear modelo de casos de uso en UML
\item Desarrollar prototipo UC funcional
\item Establecer matriz de trazabilidad Requisitos-UC-Clases
\end{itemize}

\vspace{0.5cm}

\begin{tcolorbox}[colback=yellow!5, colframe=yellow!50]
\textbf{Enfoque:} Continuar en fase de análisis sin mencionar tecnologías específicas hasta Semana 13
\end{tcolorbox}

\end{frame}

\begin{frame}
\frametitle{¡Preguntas?}

\begin{center}
\Large ¡Gracias por su atención! \\

\vspace{1cm}

\textbf{¿Preguntas?}

\vspace{1cm}

\normalsize
\textbf{Grupo 6}\\
Universidad Nacional de Ingeniería\\
CC341 - Ingeniería de Software\\
Ciclo 2025-2

\end{center}

\end{frame}

\end{document}

